\section{Ethics, data management and dissemination}
The study is prospective and requires institutional ethics review, an approved data-management plan and documented consent before recruitment. The interviewer will not provide psychotherapy, request detailed traumatic disclosures or diagnose participants. If distress occurs, the interviewer will pause, offer the option to skip or terminate, and follow an institutionally approved distress and immediate-risk procedure. Participants will receive relevant support information; confidentiality limits and the escalation procedure must be communicated before consent. If audio recording is declined, proceed only with an institutionally approved note-based option and explicit consent; otherwise do not conduct the interview. Distinguish such records from transcripts and assess their interpretive limitations rather than treating them as verbatim accounts. Confirm a private interview setting where feasible and explain that recordings cannot be fully anonymous. Disable unauthorised cloud recording, automated transcription and meeting assistants. If the connection fails, reconnect only using the agreed contact route; do not resume recording until consent is reconfirmed. Participation and refusal must not affect education, employment or access to services.

Recordings, transcripts and consent records will be stored only in approved institutional storage. No artefact transfer or retention is part of the default protocol. Live showing is optional and unrecorded; disable or pause capture and transcription, avoid copying links or excerpts, and reconfirm consent before resuming. Any exceptional transfer or retention needs a separately approved procedure, consent, rights assessment, defined purpose, access and deletion rules. Separate access-controlled consent and identity records from the pseudonymised analytic corpus. Restrict access to authorised team members; do not upload recordings, verbatim transcripts, personal prompts or artefacts to external generative systems. Do not scrape or analyse creator-community posts under this protocol. Record retention duration, deletion, withdrawal deadlines and access arrangements in the approved information sheet before recruitment. Pseudonymisation is not anonymisation; distinctive stories, lyrics, voices or platform profiles can enable re-identification. Quotations require contextual disclosure review, and potentially identifying material will be omitted or safely paraphrased with clear labelling. Translations will be checked against the source meaning and not presented as verbatim Polish quotations.

Optional artefacts may contain third-party lyrics, voices or personal material. The participant's permission alone does not establish rights for public reproduction. No identifiable audio, lyrics or images will be published without appropriate rights and explicit approval. Findings will be reported using contextualised, de-identified extracts where permitted. Public materials may include the approved interview guide, analysis procedure and non-identifying aggregate summaries; raw data will not be made publicly available. Any future sharing arrangement requires institutional approval and compliance with the consent terms. Before submission, report actual approval identifier, dates, researcher roles, recruitment, deviations, funding, conflicts and author contributions; do not infer them from this protocol.
